\section{Experiment 8A: paper-readiness and contribution audit}
\label{sec:experiment8a}

\paragraph{Purpose.}
Experiment 8A is a scope and evidence audit rather than another performance
benchmark. It asks which claims survive the complete sequence through 7C, what
the paper can be about without overstating federation, and which checks remain
necessary before freezing a manuscript. The audit treats failed gates as
evidence and does not redefine success criteria after seeing results. Its frozen
record is \path{code/config/experiment_8a.json}; the machine-readable claim map
is \path{results/tables/experiment_8a_claim_audit.csv}.

\paragraph{Decision.}
The project is \emph{conditionally ready as an LPV identification and control
paper}, but it is not ready as a federated-algorithm paper. The proposed central
claim is:
\begin{quote}
Measurable operating variation and persistent fleet heterogeneity are distinct
design axes: scheduling addresses the former, whereas model specialization
addresses the latter. Their combination is needed only when both effects are
large enough relative to the required control accuracy.
\end{quote}
This is an empirical design rule within the tested benchmark, not a universal
theorem. The paper should use ``fleet'' or ``multi-system'' rather than
``federated'' in its title unless a later experiment supplies an actual
distributed method and a separately justified benefit.

\begin{table}[htbp]
\centering
\caption{Experiment 8A claim audit. ``Boundary'' denotes a supported negative
or limiting result rather than a positive method contribution.}
\label{tab:experiment8a-claims}
\begin{tabular}{p{0.08\linewidth}p{0.44\linewidth}p{0.18\linewidth}p{0.16\linewidth}}
\toprule
ID & Candidate claim & Status & Evidence\\
\midrule
C1 & Scheduling and specialization address different variation axes & Primary, supported & 2, 3, 4G, 6A, 6B\\
C2 & Structured family LPV identification enables redesign on nonlinear-plant data & Supported & 4F, 4G, 6B\\
C3 & Scheduled M4 uses fewer stored gains than the matching tested LTI grid & Supported with scope & 5A\\
C4 & Family sharing beats Global but need not beat accurate Local identification & Boundary, supported & 7A--7C\\
C5 & A federated LPV algorithm improves accuracy or control & Unsupported & None\\
\bottomrule
\end{tabular}
\end{table}

\paragraph{Defensible contributions.}
The first contribution is the two-axis regime map. Experiment 6A varies speed
envelope and family separation independently and recovers the four expected
architectures: Global LTI, Global LPV, Family LTI, and Family LPV. Experiment 6B
preserves the corner decisions under corrected varying-speed dynamics and slower
maneuvers. The second contribution is methodological: unconstrained matrix
regression on nonlinear-plant records can yield unsuitable control redesigns,
whereas the structured positive parameter-ratio estimator restores the intended
family LPV advantage and transfers across independent fleets and trajectories
in 4F--4G. The third is a bounded implementation result: in the tested search,
M4 stores 15 gain vectors while the first gridded family-LTI alternative meeting
the 5\% tracking tolerance stores 30. The fourth is a boundary result:
Experiments 7A--7C consistently show that oracle-family sharing reduces error
relative to Global pooling but cannot outperform a highly data-efficient Local
structured fit.

\paragraph{Claims explicitly excluded.}
The present evidence does not support federated optimization, learned clustering,
privacy preservation, communication efficiency, online adaptation, formal
closed-loop stability, a universal model-selection law, or general superiority
of LPV control. Frozen-radius audits and amplitude feasibility are empirical
checks, not stability proofs. The complexity comparison is conditional on the
tested gain grids and switching rule. Results use synthetic vehicle families,
prescribed longitudinal speed, pure lateral tire forces, known scheduling, and
measured sideslip proxy and yaw rate.

\paragraph{Minimum work before manuscript freeze.}
Only three items are classified as required. First, independently regenerate
the headline numbers and figures from saved configurations and verify that every
paper value traces to an artifact. Second, add one targeted higher-fidelity
robustness check that introduces a genuinely absent coupling, preferably
longitudinal dynamics with combined-slip tire forces; this should test the
regime-map conclusion rather than launch a new method search. Third, sharpen the
baseline and literature comparison against established local-LPV and multi-model
identification/control approaches. Real vehicle data, learned family assignment,
formal robust synthesis, and communication studies would strengthen the work but
are not prerequisites for the scoped simulation paper.

\paragraph{Candidate manuscript blueprint.}
The recommended title is \emph{Scheduling or Specialization? Structured LPV
Identification and Control under Operating and Fleet Heterogeneity}. After the
introduction, the paper should define the two variation axes and four candidate
architectures, present the nonlinear benchmark and structured estimator, report
the regime map and robustness result, then give the redesign and complexity
evidence. A compact final subsection should report 7A--7C as the boundary of
collaboration: model sharing is valuable relative to a global model and for
model-count reduction, but does not improve accuracy over Local here. The full
diagnostic path, including failed 4D and 4D-R designs, belongs in supplementary
material rather than the main narrative.

\paragraph{Decision on an FL extension.}
The current manuscript should not be delayed by implementing distributed
optimization of the failed Family objective. Any FL extension starts with a new
oracle mechanism and an application-grounded benefit. It enters this paper only
if that gate adds a result that materially changes the contribution; otherwise
it should become a separate study. This separation protects both stories: the
LPV paper remains evidence-complete, while a future FL method is judged against
Local rather than only against an intentionally weak Global baseline.
